\begin{lemma}[Uniform subquadratic growth of the centered class]
Let $F$ be a finite-dimensional real inner-product space, let
$1\le C_0<\infty$, and let $z_c$ be either $(0,0)$ or $(0,1)$. Then, for every
compact $J\subseteq\R$,
\begin{equation}\label{eq:uniform-subquadratic-growth}
  \lim_{M\to\infty}\
  \sup_{\substack{v\in\Aclass_{m,C_0}\text{ centered at }z_c\\
                  |x|\ge M,\ t\in J}}
       \frac{|v(x,t)|}{|x|^2}=0 ,
\end{equation}
where $\Aclass_{m,C_0}$ is the class of \noderef{IsFitClass} and ``centered at
$z_c$'' is meant with constant $C_0$ in the sense of \noderef{IsCenteredAt}.
Explicitly: for every compact $J\subseteq\R$ and every $\varepsilon>0$ there is
$M>0$ such that $|v(x,t)|\le\varepsilon|x|^2$ for every such $v$ and whenever
$|x|\ge M$ and $t\in J$.
\end{lemma}
\begin{proof}
Let $C_1>0$ be the constant of \noderef{CenteredGrowth} for this $C_0$. Let
$J\subseteq\R$ be compact and let $\varepsilon>0$.

Since $J$ is compact it is bounded, so there is $T\ge2$ with $|t|\le T$ for all
$t\in J$. For $M\ge T$ and a point $z=(x,t)$ with $|x|\ge M$ and $t\in J$, set
$K=|x|$. Then $K\ge M\ge T\ge2$, $|x|\le K$ and $|t|\le T\le K$, so
$z\in\mathcal B_K$ and \noderef{CenteredGrowth} applies to any
$v\in\Aclass_{m,C_0}$ centered at $z_c$ — its hypothesis is that $v$ be centered
at $(0,0)$ or at $(0,1)$, and $z_c$ is one of those two points by assumption, so
that disjunction holds in either case — giving
\[
  |v(x,t)|\le C_1K\bigl(1+\log K\bigr)^2 .
\]
Dividing by $|x|^2=K^2$,
\[
  \frac{|v(x,t)|}{|x|^2}\le C_1\frac{(1+\log K)^2}{K}.
\]
The function $K\mapsto C_1(1+\log K)^2/K$ tends to $0$ as $K\to\infty$. Indeed,
for $K\ge1$ we have $\log K\ge0$ and $(\log K-1)^2\ge0$, so
$2\log K\le1+(\log K)^2$ and therefore
\[
  (1+\log K)^2=1+2\log K+(\log K)^2\le2+2(\log K)^2,
  \qquad\text{whence}\qquad
  C_1\frac{(1+\log K)^2}{K}\le2C_1\Bigl(\frac1K+\frac{(\log K)^2}{K}\Bigr),
\]
and both $1/K$ and $(\log K)^2/K$ tend to $0$ as $K\to\infty$. Hence there is
$M_0\ge2$ with
$C_1(1+\log K)^2/K\le\varepsilon$ for all $K\ge M_0$. Taking
$M=\max\{T,M_0\}>0$ gives $|v(x,t)|\le\varepsilon|x|^2$ for every centered
member $v$ of the class and all $|x|\ge M$, $t\in J$. The constant $M$ depends
only on $C_0$, $\dim F$, $J$ and $\varepsilon$, not on $v$, which is the
asserted uniformity. This is the explicit form of the statement.

\emph{The displayed limit.} It remains to read the display
\eqref{eq:uniform-subquadratic-growth} off that form. This step is not
formality: the display is a statement about a supremum, and by the convention of
the preamble a supremum carries the value $0$ when its set is empty or unbounded
above, so the boundedness has to be exhibited. Write
\[
  S(M)=\Bigl\{\frac{|v(x,t)|}{|x|^2}\ :\
      v\in\Aclass_{m,C_0}\text{ centered at }z_c\text{ with constant }C_0,\
      |x|\ge M,\ t\in J\Bigr\},
\]
the set over which the supremum in \eqref{eq:uniform-subquadratic-growth} is
taken, so that the display asserts $\sup S(M)\to0$ as $M\to\infty$.

Fix $\varepsilon>0$ and write $M_\varepsilon=\max\{T,M_0\}$ for the radius
produced above for this $\varepsilon$, which the previous paragraph called $M$;
in the present paragraph the letter $M$ is free again and ranges over the
argument of the supremum. Let $M\ge M_\varepsilon$. Every element of $S(M)$
is of the form $|v(x,t)|/|x|^2$ with $v$ centered, $t\in J$ and
$|x|\ge M\ge M_\varepsilon$, so the bound of the previous paragraph applies at
that very point and gives $|v(x,t)|\le\varepsilon|x|^2$; since $|x|\ge M>0$ we
may divide, and $|v(x,t)|/|x|^2\le\varepsilon$. Also
$|v(x,t)|/|x|^2\ge0$. Hence
\[
  S(M)\subseteq[0,\varepsilon]\qquad(M\ge M_\varepsilon).
\]
In particular $S(M)$ is bounded above, so $\sup S(M)$ is a genuine least upper
bound rather than the convention value whenever $S(M)\ne\emptyset$, and then
$0\le\sup S(M)\le\varepsilon$; and if $S(M)=\emptyset$ then $\sup S(M)=0$ by the
convention, so $0\le\sup S(M)\le\varepsilon$ in that case too. Thus
$0\le\sup S(M)\le\varepsilon$ for every $M\ge M_\varepsilon$. As $\varepsilon>0$
was arbitrary, $\sup S(M)\to0$ as $M\to\infty$, which is
\eqref{eq:uniform-subquadratic-growth}.
\end{proof}
